% page 395 of the facsimile (image p-394 of the scan)
% block 175.3 x 237.7 mm ; left margin 26.7 mm
\pgc{5.080mm}{0.000mm}{
\l{\cel{55}387}
\l{}
\l{\cel{4}\sou{musho}. (\sou{zool.}) Insekto diptera, ek la familio \textquotedbl{}muskidi\textquotedbl{}. -}
\l{\cel{7}L.\cel{1}\sou{musca}.\cel{2}FIRS.}
\l{}
\l{\cel{4}\sou{musko}. (\sou{bot.)}\cel{2}Planto kriptogama celuloza, di qua la folieti}
\l{\cel{7}tenua multigas su tale ke li tapetizas la loki ube lu kreskas.}
\l{\cel{7}L.\cel{1}\sou{muscua.}\cel{2}FI.}
\l{}
\l{\cel{4}\sou{muskado}. (\sou{bot.}) Nuco aromatoza, quan produktas arboro exotika,}
\l{\cel{7}uzata kom spico. - L.\cel{1}\sou{myristica fragrans.}\cel{2}- DFIRS.}
\l{}
\l{\cel{4}\sou{muskato}. (\sou{bot.}) Vitbero di qua la odoro similesas olta di musko.}
\l{\cel{7}- DEFIRS.}
\l{}
\l{\cel{4}\sou{musketono(olim).}\cel{2}Pafarmo plu grosa, ma min longa, kam musketo. -}
\l{\cel{7}II. Speco di fusilo kurta, inter karabino e pistolo. -}
\l{}
\l{\cel{4}\sou{musketo}.\cel{2}(\sou{olim}). Pafarmo quan onu funcionigis per mecho acen-}
\l{\cel{7}dita e qua remplasesis da fusilo. - DEFIRS.}
\l{}
\l{\cel{4}\sou{muskulo}. (\sou{anat}.) Organo karnoza, ek fibri kontraktebla, qui}
\l{\cel{7}produktasla movi (che la animali). - DEFIRS.}
\l{}
\l{\cel{4}\sou{muslo}. (\sou{zool.}) Molusko bivalva, manjebla, tipo di la tribuo}
\l{\cel{7}\textquotedbl{}mitilinei\textquotedbl{}. - L. mytilus edulis. - DEF.}
\l{}
\l{\cel{4}muslino. Stofo de kotono tenua, e di qua la texuro esas diafana.}
\l{\cel{7}- DEFIRS.}
\l{}
\l{\cel{4}musono. (meteorol.) Vento qua aparas periodale en la oceano}
\l{\cel{7}di India, de esto a westo, en la semestro unesma di la yaro,}
\l{\cel{7}ed inversa-sinse en la semestro duesma. - DEFIRS.}
\l{}
\l{\cel{4}\sou{mustar}. (\sou{trans.)}\cel{2}Koaktesar ne-eskapeble, ne-eviteble. -\cel{2}DE.}
\l{}
\l{\cel{4}mustardo. Kondimento ek grani di sinapo, triturita kun vino-mosto}
\l{\cel{7}o vinagro. - EFIS.}
\l{}
\l{\cel{4}\sou{muta}. Qua ne povas (fiziologiale) parolar. - EFIS.}
\l{}
\l{\cel{4}\sou{mutaco.}\cel{2}(\sou{biol.)}\cel{2}Vario bruska e spontana qua eventas en decen}
\l{\cel{7}danto di individuo normala, e qua, de lore, esas fixa e here-}
\l{\cel{7}donta. - DEFIS.}
\l{}
\l{\cel{4}\sou{mutilar}. (\sou{trans.)}\cel{2}I. Alterar (la korpo) ye sua integreso per}
\l{\cel{7}deprenar de lu membro od altra parto. -II. (\sou{metaf.}) Alterar}
\l{\cel{7}(ensemblo) ye lua integreso, per deprenar parto esencala. -}
\l{\cel{7}EFIS.}
\l{}
\l{\cel{4}\sou{mutono}. (\sou{zool.}) Mamifero ek la familio \textquotedbl{}rumineri\textquotedbl{},\cel{1}\sou{kun}\cel{1}chanfreno}
\l{\cel{7}konvexa, korni kava e permananta, pili lana, frizita. L. ovis}
\l{\cel{7}\sou{aries}.\cel{2}II. (\sou{meteorol.}) Mikra nubi, ek ondi blanka, de vaporo,}
\l{\cel{7}simila a lano-floki. - III. (\sou{relig.katol.}) Persono konfidita}
\l{\cel{7}a la direkto de pastoro spiritala. - EFI.}
}
